\chapter{Thực nghiệm và đánh giá}

\section{Bộ dữ liệu}
RACE (Reading Comprehension Dataset From Examinations) là một bộ dữ liệu quy mô lớn trong lĩnh vực xử lý ngôn ngữ tự nhiên. Khác với bộ dữ liệu SQuAD (chủ yếu dựa trên Wikipedia), RACE được xây dựng trực tiếp từ đề thi của các kỳ thi tiếng Anh thực tế dành cho học sinh trung học cơ sở và trung học phổ thông tại hệ thống giáo dục Trung Quốc.

Hai lý do chính nghiên cứu sử dụng RACE làm bộ dữ liệu đánh giá mô hình là:
Độ phức tạp và đa dạng ngữ nghĩa: Văn bản trong RACE thường bao trùm nhiều lĩnh vực trong đời sống: giáo dục, tin tức, văn hóa, chính trị, khoa học và các câu chuyện luân lý. Chính vì thế, đặc điểm dữ liệu và văn phong của văn bản trong RACE vô cùng sát với các tài liệu trong thực tế. Các văn bản trong RACE có độ dài tiêu chuẩn từ 300 đến 500 từ, độ dài này đủ lớn để chứa đựng lượng thông tin cho một bài kiểm tra: các chuỗi sự kiện phức tạp, luận điểm, nhưng vẫn đảm bảo không làm quá tải cửa sổ ngữ cảnh (context window) của các sLLM.

Chiều sâu logic và tính đa nghĩa của văn bản gốc: Không giống các văn bản bách khoa toàn thư chỉ liệt kê dữ liệu bề mặt, các văn bản của RACE được các chuyên gia giáo dục tuyển chọn chuyên biệt nhằm mục đích đánh giá nhận thức của người học. Cấu trúc của văn bản ẩn chứa nhiều tầng nghĩa, tích hợp các chuỗi nhân quả, thái độ ngầm định của tác giả và các mệnh đề gây nhiễu. Nếu nghiên cứu lựa chọn một bộ dữ liệu văn bản thiếu chiều sâu, các nỗ lực đánh giá LLM sinh câu hỏi tư duy cao đều sẽ thất bại.

\section{Thiết lập thực nghiệm}
\subsection{Các mô hình sử dụng}
\subsection{Cấu hình thực nghiệm}

\section{Kết quả sinh câu hỏi}

\section{Kết quả OCR}

\section{Kết quả đánh giá định dạng đầu ra}

\section{Thảo luận}

\section{Hạn chế và hướng phát triển}